\SourceLeaf{102}
\textit{Demonstr.} I. Quando $p=2$, in (I) omnes termini praeter primum, \textit{i.\ e.} $m$ termini divisibiles erunt; totidem autem erunt in (II).

II. Sit $p$ numerus impar vel numeri imparis duplum vel quadruplum, atque $a\equiv rr\ (\mathrm{mod.}\,p)$. Tum in progressione, $-m,-(m-1),-(m-2),.....+m$ (quae terminorum multitudine cum $(\mathrm{II})$ convenit et per (III) designabitur) totidem ad minimum termini erunt secundum modulum $p$ ipsi $r$ congrui, quot in serie (II) per $p$ divisibiles (art. praec.). Inter illos autem bini, qui signo tantum, non magnitudine, discrepent, occurrere nequeunt\HistoricalNote{*)}{Si enim esset $r\equiv-f\equiv+f\,(\operatorname{mod.}p)$, fieret $2f$ per $p$ divisibilis, adeoque etiam $2a$ (propter $ff\equiv a\,(\operatorname{mod.}p)$). Hoc autem aliter fieri nequit, quam si $p=2$, quum per hyp. $a$ ad $p$ sit primus. Sed de hoc casu iam seorsim diximus.}. Tandem quisque eorum correspondentem habebit in serie (I), qui per $p$ erit divisibilis. Scilicet si fuerit $\pm b$ aliquis terminus seriei (III) ipsi $r$ secundum $p$ congruus, erit $a-bb$ per $p$ divisibilis. Quodsi igitur $b$ est par, terminus seriei (I), $2(a-bb)$, per $p$ divisibilis erit. Si vero $b$ impar, terminus $\frac12(a-bb)$ per $p$ divisibilis erit: namque manifesto $\frac{a-bb}{p}$ erit integer \textit{par}, quoniam $a-bb$ per $8$, $p$ autem ad summum per $4$ divisibilis ($a$ enim per hyp. est formae $8n+1$, $bb$ autem ideo quod est numeri imparis quadratum eiusdem formae erit, quare differentia erit formae $8n$). Hinc tandem concluditur, in serie (I) totidem terminos esse per $p$ divisibiles, quot in (III) sint ipsi $r$ secundum $p$ congrui \textit{i.\ e.} totidem aut plures quam in (II) sint per $p$ divisibiles. \textit{Q.\ E.\ D.}

III. Sit $p$ formae $8n$, atque $a\equiv rr\ (\mathrm{mod.}\,2p)$. Facile enim perspicitur, $a$, quum ex hyp. ipsius $p$ sit residuum, etiam ipsius $2p$ residuum fore. Tum in serie (III) totidem ad minimum termini erunt ipsi $r$ secundum $p$ congrui, quot in (II) sunt per $p$ divisibiles, illique omnes magnitudine erunt inaequales. At cuique eorum respondebit aliquis in (I) per $p$ divisibilis. Si enim $+b$ vel $-b\equiv r\ (\mathrm{mod.}\,p)$, erit $bb\equiv rr\ (\mathrm{mod.}\,2p)$\HistoricalNote{\(\dagger\))}{Erit scilicet $bb-rr=(b-r)(b+r)$ e duobus factoribus compositus, quorum alter per $p$ divisibilis (hyp.), alter per $2$ (quia tum $b$ tum $r$ sunt impares); adeoque $bb-rr$ per $2p$ divisibilis.}, adeoque terminus $\frac12(a-bb)$ per $p$ divisibilis. Quare in (I) totidem ad minimum termini erunt per $p$ divisibiles quam in (II). \textit{Q.\ E.\ D.}

\Article{129}

\textsc{Theorema.} \textit{Si $a$ est numerus primus formae $8n+1$, necessario infra $2\surd a+1$ dabitur aliquis numerus primus cuius non-residuum sit $a$.}

\SourceLeaf{103}
\textit{Demonstr.} Esto, si fieri potest, $a$ residuum omnium primorum ipso $2\surd a+1$ minorum. Tum facile perspicietur, $a$ etiam omnium numerorum compositorum ipso $2\surd a+1$ minorum residuum fore (conferantur praecepta per quae diiudicare docuimus, utrum numerus propositus sit numeri compositi residuum necne; art. $105$). Sit numerus proxime minor quam $\surd a$, $=m$. Tum in serie

\[
\begin{gathered}(\mathrm I)......a,\tfrac12(a-1),2(a-4),\tfrac12(a-9)....2(a-mm)\\ \text{vel}\quad\tfrac12(a-mm)\end{gathered}
\]

totidem aut plures termini erunt per numerum quemcunque ipso $2\surd a+1$ minorem divisibiles, quam in hac

\[
(\mathrm{II}).....1,2,3,4....2m+1\quad\text{(art. praec.)}
\]

Hinc vero sequitur, productum ex omnibus terminis (I) per productum omnium terminorum (II) divisibile esse, (art. $126$). At illud est aut \mbox{$=a(a-1)(a-4)....(a-mm)$} aut semissis huius producti (prout $m$ aut par aut impar). Quare productum $a(a-1)(a-4)....(a-mm)$ certo per productum omnium terminorum $(\mathrm{II})$ dividi poterit, et, quia omnes hi termini ad $a$ sunt primi, etiam productum illud omisso factore $a$. Sed productum ex omnibus terminis $(\mathrm{II})$ ita etiam exhiberi potest,

\[
(m+1).((m+1)^2-1).((m+1)^2-4)....((m+1)^2-m^2)
\]

Fiet igitur

\[
\frac1{m+1}.\frac{a-1}{(m+1)^2-1}.\frac{a-4}{(m+1)^2-4}....\frac{a-m^2}{(m+1)^2-m^2}
\]

numerus integer, quamquam sit productum ex fractionibus unitate minoribus: quia enim necessario $\surd a$ irrationalis esse debet, erit $m+1>\surd a$, adeoque $(m+1)^2>a.$ Hinc tandem concluditur suppositionem nostram locum habere non posse. \textit{Q.\ E.\ D.}

Iam quia $a$ certo $>9$, erit $2\surd a+1<a$, dabiturque adeo aliquis primus $<a$ cuius non-residuum $a$.

\begin{center}\textit{Per inductionem theorema generale (fundamentale) stabilitur, conclusionesque inde deducuntur.}\end{center}

\Article{130}

Postquam rigorose demonstravimus quemvis numerum primum formae $4n+1$, et positive et negative acceptum, alicuius numeri primi ipso minoris non-\SourceLeaf{104}residuum esse, ad comparationem exactiorem et generaliorem numerorum primorum, quatenus unus alterius residuum vel non-residuum est, statim transimus.

Omni rigore supra demonstravimus, $-3$ et $+5$ esse residua vel non-residua omnium numerorum primorum, qui ipsorum $3, 5$ respective sint residua vel non-residua.

Per inductionem autem circa numeros sequentes institutam invenitur: $-7$, $-11$, $+13$, $+17$, $-19$, $-23$, $+29$, $-31$, $+37$, $+41$, $-43$, $-47$, $+53$, $-59$ etc. esse residua vel non-residua omnium numerorum primorum, qui, positive sumti, illorum primorum respective sint residua vel non-residua. Inductio haec perfacile adiumento tabulae II confici potest.

Quivis autem levi attentione adhibita observabit, ex his numeris primis signo positivo affectos esse eos, qui sint formae $4n+1$, negativo autem eos, qui sint formae $4n+3$.

\Article{131}

Quod hic per inductionem deteximus, generaliter locum habere mox demonstrabimus. Antequam autem hoc negotium adeamus, necesse erit, omnia quae ex theoremate, si verum esse supponitur, sequuntur, eruere. Theorema ipsum ita enunciamus.

\emph{Si $p$ est numerus primus formae $4n+1$, erit $+p$, si vero $p$ formae $4n+3$, erit $-p$ residuum vel non-residuum cuiusvis numeri primi qui positive acceptus ipsius $p$ est residuum vel non-residuum.}

Quia omnia fere quae de residuis quadraticis dici possunt, huic theoremati innituntur, denominatio \emph{theorematis fundamentalis}, qua in sequentibus utemur, haud absona erit.

Ut ratiocinia nostra quam brevissime exhiberi possint, per $a, a', a''$ etc. numeros primos formae $4n+1$, per $b, b', b''$ etc. numeros primos formae $4n+3$ denotabimus; per $A, A', A''$ etc. numeros quoscunque formae $4n+1$, per $B, B', B''$ etc. autem numeros quoscunque formae $4n+3$; tandem litera $R$ duabus quantitatibus interposita indicabit, priorem sequentis esse residuum, sicuti litera $N$ significationem contrariam habebit. \emph{Ex. gr.} $+5R11$, $\pm 2N5$, indicabit $+5$ ipsius $11$ esse residuum, $+2$ vel $-2$ esse ipsius $5$ non-residuum. Iam col\SourceLeaf{105}lato theoremate fundamentali cum theorematibus art. $111$, sequentes propositiones facile deducentur.

\[
\begin{array}{rccc}&\text{Si}&&\text{erit}\\[5pt]
1.&\pm aRa'&\mathinner{.\ .\ .\ .\ .\ .\ .}&\pm a' Ra\\[4pt]
2.&\pm aNa'&\mathinner{.\ .\ .\ .\ .\ .\ .}&\pm a' Na\\[4pt]
3.&\left\{\begin{array}{l}+aRb\\-aNb\end{array}\right\}&\mathinner{.\ .\ .\ .\ .\ .\ .}&\pm bRa\\[4pt]
4.&\left\{\begin{array}{l}+aNb\\-aRb\end{array}\right\}&\mathinner{.\ .\ .\ .\ .\ .\ .}&\pm bNa\\[4pt]
5.&\pm bRa&\mathinner{.\ .\ .\ .\ .\ .\ .}&\left\{\begin{array}{l}+aRb\\-aNb\end{array}\right.\\[4pt]
6.&\pm bNa&\mathinner{.\ .\ .\ .\ .\ .\ .}&\left\{\begin{array}{l}+aNb\\-aRb\end{array}\right.\\[4pt]
7.&\left\{\begin{array}{l}+bRb'\\-bNb'\end{array}\right\}&\mathinner{.\ .\ .\ .\ .\ .\ .}&\left\{\begin{array}{l}+b' Nb\\-b' Rb\end{array}\right.\\[4pt]
8.&\left\{\begin{array}{l}+bNb'\\-bRb'\end{array}\right\}&\mathinner{.\ .\ .\ .\ .\ .\ .}&\left\{\begin{array}{l}+b' Rb\\-b' Nb\end{array}\right.\\[4pt]
\end{array}
\]

\Article{132}

In his omnes casus, qui, duos numeros primos comparando. occurrere possunt, continentur: quae sequuntur, ad numeros quoscunque pertinent: sed harum demonstrationes minus sunt obviae.

\[
\begin{array}{rccc}&\text{Si}&&\text{erit}\\[5pt]
9.&\pm aRA&\mathinner{.\ .\ .\ .\ .\ .\ .}&\pm ARa\\[4pt]
10.&\pm bRA&\mathinner{.\ .\ .\ .\ .\ .\ .}&\left\{\begin{array}{l}+ARb\\-ANb\end{array}\right.\\[4pt]
11.&+aRB&\mathinner{.\ .\ .\ .\ .\ .\ .}&\pm BRa\\[4pt]
12.&-aRB&\mathinner{.\ .\ .\ .\ .\ .\ .}&\pm BNa\\[4pt]
13.&+bRB&\mathinner{.\ .\ .\ .\ .\ .\ .}&\left\{\begin{array}{l}-BRb\\+BNb\end{array}\right.\\[4pt]
14.&-bRB&\mathinner{.\ .\ .\ \phantom{.\ .\ .}\ .}&\left\{\begin{array}{l}+BRb\\-BNb\end{array}\right.\\[4pt]
\end{array}
\]

\SourceLeaf{106}
Quum omnium harum propositionum demonstrationes ex iisdem principiis sint petendae, necesse non erit omnes evolvere: demonstratio prop. $9$, quam apponimus tamquam exemplum inservire potest. Ante omnia autem observetur, quemvis numerum formae $4n+1$ aut nullum factorem formae $4n+3$ habere, aut duos, aut quatuor etc., \emph{i. e.} multitudinem talium factorum (inter quos etiam aequales esse possunt) semper fore parem: quemvis vero formae $4n+3$ multitudinem imparem factorum formae $4n+3$ (\emph{i. e.} aut unum aut tres aut quinque etc.) implicare. Multitudo factorum formae $4n+1$ indeterminata manet.

\emph{Prop.} $9$ ita demonstratur. Sit $A$ productum e factoribus primis $a',a'',a'''$ etc., $b,b',b''$ etc.; eritque factorum $b,b',b''$ etc. multitudo par (possunt etiam nulli adesse, quod eodem redit). Iam si $a$ est residuum ipsius $A$, erit residuum etiam omnium factorum $a',a'',a'''$ etc. $b,b',b''$ etc. quare per propp. $1,3$ art. praec. singuli hi factores erunt residua ipsius $a$, adeoque etiam productum $A$. $-A$ vero idem esse debet. --- Quodsi vero $-a$ est residuum ipsius $A$, eoque ipso omnium factorum $a',a''$ etc. $b,b'$ etc.; singuli $a',a''$ etc. erunt ipsius $a$ residua, singuli $b,b'$ etc. autem non-residua. Sed quum posteriorum multitudo sit par, productum ex omnibus, \emph{i. e.} $A$, ipsius $a$ residuum erit, hincque etiam $-A$.

\Article{133}

Investigationem adhuc generalius instituamus. Contemplemur duos numeros quoscunque impares inter se primos, signis quibuscunque affectos, $P$ et $Q$. Concipiatur $P$ sine respectu signi sui in factores suos primos resolutus, designeturque per $p$, quot inter hos reperiantur quorum non-residuum sit $Q$. Si vero aliquis numerus primus, cuius non-residuum est $Q$, pluries inter factores ipsius $P$ occurrit, pluries etiam numerandus erit. Similiter sit $q$ multitudo factorum primorum ipsius $Q$, quorum non-residuum est $P$. Tum numeri $p,q$ certam relationem mutuam habebunt ab indole numerorum $P,Q$ pendentem. Scilicet si alter numerorum $p,q$ est par vel impar, numerorum $P,Q$ forma docebit, utrum alter par sit vel impar. Haec relatio in sequenti tabula exhibetur.

Erunt $p,q$ simul pares vel simul impares, quando numeri $P,Q$ habent formas:

\[
\begin{array}{rl}1.&+A,\quad+A'\\2.&+A,\quad-A'\\\SourceLeaf{107}3.&+A,\quad+B\\4.&+A,\quad-B\\5.&-A,\quad-A'\\6.&+B,\quad-B'\end{array}
\]

Contra numerorum $p,q$ alter erit par, alter impar, quando numeri $P,Q$ habent formas:

\[
\begin{array}{rl}7.&-A,\quad+B\\8.&-A,\quad-B\\9.&+B,\quad+B'\\10.&-B,\quad-B'{}^{*)}\end{array}
\]

\HistoricalDisplayNote{*)}{\begin{tabular}{@{}ll@{}}Sit &$l=1$ si uterque $P,Q\equiv3\,(\operatorname{mod.}4)$, alioquin $l=0$\\&$m=1$ si uterque $P,Q$ negativus, alioquin $m=0$\end{tabular}\par tunc relatio pendet ab $l+m$.}

\emph{Ex.} Sint numeri propositi $-55$ et $+1197$, qui ad casum quartum erunt referendi. Est autem $1197$ non-residuum unius factoris primi ipsius $55$. scilicet numeri $5$, $-55$ autem non-residuum trium factorum primorum ipsius $1197$, scilicet numerorum $3,3,19$.

Si $P$ et $Q$ numeros primos designant, propositiones hae abeunt in eas quas art. $131$ tradidimus. Hic scilicet $p$ et $q$ maiores quam $1$ fieri nequeunt, quare quando $p$ ponitur esse par necessario erit $=0$ \emph{i. e.} $Q$ erit residuum ipsius $P$, quando vero $p$ est impar, $Q$ ipsius $P$ non-residuum erit. Et vice versa. Ita scriptis $a,b$ loco ipsorum $A,B$, ex $8$ sequitur, si $-a$ fuerit residuum vel non-residuum ipsius $b$, fore $-b$ non-residuum vel residuum ipsius $a$, quod cum $3$ et $4$ art. $131$ convenit.

Generaliter vero patet, $Q$ residuum ipsius $P$ esse non posse nisi fuerit $p=0$; si igitur $p$ impar, $Q$ certo ipsius $P$ non-residuum erit.

Hinc etiam propp. art. praec. sine difficultate derivari possunt.

Ceterum mox patebit, hanc repraesentationem generalem plus esse quam speculationem sterilem, quum theorematis fundamentalis demonstratio completa absque ea vix perfici possit.

\SourceLeaf{108}
\Article{134}

Aggrediamur nunc deductionem harum propositionum.

I.\quad Concipiatur, ut ante, $P$ in factores suos primos resolutus, signis neglectis, insuperque etiam $Q$ in factores quomodocunque resolvatur, ita tamen ut signi ipsius $Q$ ratio habeatur. Combinentur illi singuli cum singulis his. Tum si $s$ designat multitudinem omnium combinationum, in quibus factor ipsius $Q$ est non-residuum factoris ipsius $P$, $p$ et $s$ vel simul pares vel simul impares erunt. Sint enim factores primi ipsius $P$, hi $f,f',f''$ etc. et inter factores in quibus $Q$ est resolutus, sint $m$ qui ipsius $f$ sint non-residua, $m'$ non-residua ipsius $f'$ $m''$ non-residua ipsius $f''$ etc. Tum facile quisquis perspiciet, fore

\[
s=m+m'+m''+\text{etc.}
\]

$p$ autem exprimere quot numeri inter ipsos $m,m',m''$ etc. sint impares. Unde sponte patet, $s$ fore parem quando $p$ sit par, imparem quando $p$ sit impar.

II.\quad Haec generaliter valent, quomodocunque $Q$ in factores sit resolutus. Descendamus ad casus particulares. Contemplemur primo casus, ubi alter numerorum, $P$, est positivus, alter vero, $Q$, vel formae $+A$ vel formae $-B$. Resolvantur $P,Q$ in factores suos primos, attribuatur singulis factoribus ipsius $P$ signum positivum, singulis autem factoribus ipsius $Q$ signum positivum vel negativum, prout sunt formae $a$ vel $b$; tunc autem manifesto $Q$ fiet vel formae $+A$ vel $-B$ uti requiritur. Combinentur factores singuli ipsius $P$ cum singulis factoribus ipsius $Q$, designetque ut ante $s$ multitudinem combinationum in quibus factor ipsius $Q$ est non-residuum factoris ipsius $P$, similiterque $t$ multitudinem combinationum in quibus factor ipsius $P$ est non-residuum factoris ipsius $Q$. At ex theoremate fundamentali sequitur illas combinationes identicas fore cum his adeoque $s=t$. Tandem ex iis quae modo demonstravimus sequitur esse $p\equiv s\, (\mathrm{mod.}\ 2)$, $q\equiv t\, (\mathrm{mod.}\ 2)$, unde fit $p\equiv q\, (\mathrm{mod.}\ 2)$.

Habentur itaque propp. $1,3,4$ et $6$ art. $133$.

Propositiones reliquae per methodum similem directe erui possunt, sed una consideratione nova indigent; facilius autem ex praecedentibus sequenti modo derivantur.

III.\quad Denotent rursus $P,Q$, numeros quoscunque impares inter se primos, $p,q$ multitudinem factorum primorum ipsorum $P,Q$, quorum non-residua $Q$, $P$ respective. Tandem sit $p'$ multitudo factorum primorum ipsius $P$, quorum \SourceLeaf{109}non-residuum est $-Q$ (quando $Q$ per se est negativus, manifesto $-Q$ numerum positivum indicabit). Iam omnes factores primi ipsius $P$ in quatuor classes distribuantur.

1) in factores formae $a$, quorum residuum est $Q$.

2) factores formae $b$, quorum residuum $Q$. Horum multitudo sit $\chi$.

3) factores formae $a$, quorum non-residuum est $Q$. Horum multitudo sit $\psi$.

4) factores formae $b$, quorum non-residuum $Q$. Quorum multitudo $=\omega$.

Tum facile perspicitur fore $p=\psi+\omega$, $p'=\chi+\psi$.

Iam quando $P$ est formae $\pm A$, erit $\chi+\omega$ adeoque etiam $\chi-\omega$ numerus par: quare fiet $p'=p+\chi-\omega\equiv p\, (\mathrm{mod.}\ 2)$; quando vero $P$ est formae $\pm B$, per simile ratiocinium invenitur, numeros $p,p'$ sec. mod. $2$ incongruos fore.

IV.\quad Applicemus haec ad casus singulos. Sit primo tum $P$ tum $Q$ formae $+A$, eritque ex prop. $1$ $p\equiv q\, (\mathrm{mod.}\ 2)$; at erit $p'\equiv p\, (\mathrm{mod.}\ 2)$; quare etiam $p'\equiv q\ (\mathrm{mod.}\ 2)$. Quod convenit cum prop. $2$. --- Simili modo si $P$ est formae $-A$. $Q$ formae $+A$, erit $p\equiv q\, (\mathrm{mod.}\ 2)$ ex prop. $2$ quam modo demonstravimus; hinc. ob $p'\equiv p$, erit $p'\equiv q$. Est itaque etiam prop. $5$ demonstrata.

Eodem modo prop. $7$ ex $3$; prop. $8$ vel ex $4$ vel ex $7$; prop. $9$ ex $6$; ex eademque prop. $10$ derivantur.

\begin{center}\emph{Demonstratio rigorosa theorematis fundamentalis.}\end{center}

\Article{135}

Per art. praec. propositiones art. $133$ non quidem sunt demonstratae, sed tamen earum veritas a veritate theorematis fundamentalis quam aliquantisper supposuimus pendere ostensa est. At ex ipsa deductionis methodo manifestum est. illas valere pro numeris $P,Q$, si modo theorema fundamentale pro omnibus factoribus primis horum numerorum inter se comparatis locum habeat, etiamsi generaliter verum non sit. Nunc igitur ipsius theorematis fundamentalis demonstrationem aggrediamur. Cui praemittimus sequentem explicationem.

\emph{Theorema fundamentale usque ad numerum aliquem $M$ verum esse dicemus, si valet pro duobus numeris primis quibuscunque, quorum neuter ipsum $M$ superat.}

Simili modo intelligi debet, si theoremata artt. $131,132,133$ usque ad aliquem terminum vera esse dicemus. Facile vero perspicitur, si de veritate theorematis fundamentalis usque ad aliquem terminum constet. has propositiones usque ad eundem terminum locum esse habituras.

\SourceLeaf{110}
\Article{136}

Theorema fundamentale pro numeris parvis verum esse, per inductionem facile confirmari, atque sic limes determinari potest usque ad quem certo locum teneat. Hanc inductionem institutam esse postulamus: prorsus autem indifferens est quousque eam persecuti simus; sufficeret adeo, si tantummodo usque ad numerum $5$ eam confirmavissemus, hoc autem per unicam observationem absolvitur, quod est $+5N3$, $\pm3N5$.

Iam si theorema fundamentale generaliter verum non est, dabitur limes aliquis, $T$, vsque ad quem valebit, ita tamen ut usque ad numerum proxime maiorem, $T+1$, non amplius valeat. Hoc autem idem est ac si dicamus, dari duos numeros primos quorum maior sit $T+1$, et qui inter se comparati theoremati fundamentali repugnent, binos autem alios numeros primos quoscunque, si modo ambo ipso $T+1$ sint minores, huic theoremati esse consentaneos. Unde sequitur, propositiones artt. $131,132,133$ usque ad $T$ etiam locum habituras. Hanc vero suppositionem consistere non posse nunc ostendemus. Erunt autem secundum formas diversas, quas tum $T+1$, tum numerus primus ipso $T+1$ minor, quem cum $T+1$ comparatum theoremati repugnare supposuimus. habere possunt, casus sequentes distinguendi. Numerum istum primum per $p$ designamus.

Quando tum $T+1$ tum $p$ sunt formae $4n+1$, theorema fundamentale duobus modis falsum esse posset, scilicet si simul esset, \emph{vel}

\[
\begin{array}{ll}&\pm pR(T+1)\quad\text{et}\quad\pm(T+1)Np\\\textit{vel simul}&\pm pN(T+1)\quad\text{et}\quad\pm(T+1)Rp\end{array}
\]

Quando tum $T+1$ tum $p$ sunt formae $4n+3$, theor. fund. falsum erit. si simul fuerit \emph{vel}

\[
\begin{array}{ll}&+pR(T+1)\quad\text{et}\quad-(T+1)Np\\\text{(sive quod eodem redit}&-pN(T+1)\quad\text{et}\quad+(T+1)Rp)\\\textit{vel}&+pN(T+1)\quad\text{et}\quad-(T+1)Rp\\\text{(sive}&-pR(T+1)\quad\text{et}\quad+(T+1)Np)\end{array}
\]

Quando $T+1$ est formae $4n+1$, $p$ vero formae $4n+3$; theor. fund. falsum erit, si fuerit \emph{vel}

\[
\begin{array}{ll}&\pm pR(T+1)\ \text{et}\ +(T+1)Np\ (\text{sive}\ -(T+1)Rp)\\\textit{vel}&\pm pN(T+1)\ \text{et}\ -(T+1)Np\ (\text{sive}\ +(T+1)Rp)\end{array}
\]

\SourceLeaf{111}
Quando $T+1$ est formae $4n+3$, $p$ vero formae $4n+1$, theor. fund. falsum erit, si fuerit \emph{vel}

\[
\begin{array}{ll}&+pR(T+1)\ (\text{sive}\ -pN(T+1))\ \text{et}\ \pm(T+1)Np\\\textit{vel}&+pN(T+1)\ (\text{sive}\ -pR(T+1))\ \text{et}\ \pm(T+1)Rp\end{array}
\]

Si demonstrari poterit, nullum horum octo casuum locum habere posse, simul certum erit, theorematis fundamentalis veritatem nullis limitibus circumscriptam esse. Hoc itaque negotium nunc aggredimur: at quoniam alii horum casuum ab aliis sunt dependentes, eundem ordinem, quo eos hic enumeravimus, servare non licebit.

\Article{137}

\emph{Casus primus. Quando $T+1$ est formae $4n+1$ $(=a)$, atque $p$ eiusdem formae; insuper vero $\pm pRa$, non potest esse $\pm aNp$.} Hic casus supra fuit primus.

Sit $+p\equiv e^2(\mathrm{mod.}\ a)$, atque $e$ par et $<a$ (quod semper obtineri potest). Iam duo casus sunt distinguendi.

I.\quad Quando $e$ per $p$ non est divisibilis. Ponatur $e^2=p+af$, eritque $f$ positivus, formae $4n+3$ (sive formae $B$), $<a$, et per $p$ non divisibilis. Porro erit $e^2\equiv p(\mathrm{mod.}\ f)$, \emph{i. e.} $pRf$ adeoque ex prop. $11$ art. $132$ $\pm fRp$ (quia enim $p,f<a$, pro his propositiones istae valebunt). At est etiam $afRp$, quare fiet quoque $\pm aRp$.

II.\quad Quando $e$ per $p$ est divisibilis, ponatur $e=gp$, atque $e^2=p+aph$, sive $pg^2=1+ah$. Tum erit $h$ formae $4n+3$ ($B$), atque ad $p$ et $g^2$ primus. Porro erit $pg^2Rh$, adeoque etiam $pRh$, hinc (prop. $11$ art. $132$) $\pm hRp$. At est etiam $-ahRp$, quia $-ah\equiv1(\mathrm{mod.}\ p)$; quare fiet etiam $\mp aRp$.

\Article{138}

\emph{Casus secundus. Quando $T+1$ est formae $4n+1$ $(=a)$. $p$ formae $4n+3$, atque $\pm pR(T+1)$, non potest esse $+(T+1)Np$ sive $-(T+1)Rp$.} Hic casus supra fuit quintus.

Sit ut supra $e^2=p+fa$ atque $e$ par et $<a$.

I.\quad Quando $e$ per $p$ non est divisibilis, erit etiam $f$ per $p$ non divisibilis. Praeterea autem erit $f$ positivus, formae $4n+1$ (sive $A$), atque $<a$; $+pRf$,%
